\documentclass{report}

%%%%%%%%%%%%%%%%%%%%%%%%%%%%%%%%%
% PACKAGE IMPORTS
%%%%%%%%%%%%%%%%%%%%%%%%%%%%%%%%%


% Codificación y tipografía básica (se cargan aquí para orden correcto)
\usepackage[utf8]{inputenc}
\usepackage[T1]{fontenc}
\usepackage{lmodern}
% Idioma español para nombres de elementos flotantes (Figura, Tabla, etc.)
\usepackage[spanish]{babel}
% Asegurar nombre de tablas en español
\addto\captionsspanish{\renewcommand{\tablename}{Tabla}}


\usepackage[tmargin=2cm,rmargin=1in,lmargin=1in,margin=0.85in,bmargin=2cm,footskip=.2in]{geometry}
\usepackage{amsmath,amsfonts,amsthm,amssymb,mathtools}
\usepackage[varbb]{newpxmath}
\usepackage{xfrac}
\usepackage[makeroom]{cancel}
\usepackage{mathtools}
\usepackage{bookmark}
\usepackage{enumitem}
\usepackage{hyperref,theoremref}
\hypersetup{
	pdftitle={Assignment},
	colorlinks=true, linkcolor=gray,citecolor=blue,
	bookmarksnumbered=true,
	bookmarksopen=true
}
\usepackage[most,many,breakable]{tcolorbox}
\usepackage{xcolor}
\usepackage{varwidth}
\usepackage{varwidth}
\usepackage{etoolbox}
%\usepackage{authblk}
\usepackage{nameref}
\usepackage{multicol,array}
\usepackage{tikz-cd}
\usepackage[ruled,vlined,linesnumbered]{algorithm2e}
\usepackage{comment} % enables the use of multi-line comments (\ifx \fi) 
\usepackage{import}
\usepackage{xifthen}
\usepackage{pdfpages}
\usepackage{transparent}

\newcommand\mycommfont[1]{\footnotesize\ttfamily\textcolor{blue}{#1}}
\SetCommentSty{mycommfont}
\newcommand{\incfig}[1]{%
    \def\svgwidth{\columnwidth}
    \import{./figures/}{#1.pdf_tex}
}

\usepackage{tikzsymbols}
\renewcommand\qedsymbol{$\Laughey$}


%\usepackage{import}
%\usepackage{xifthen}
%\usepackage{pdfpages}
%\usepackage{transparent}


%%%%%%%%%%%%%%%%%%%%%%%%%%%%%%
% SELF MADE COLORS
%%%%%%%%%%%%%%%%%%%%%%%%%%%%%%



\definecolor{myg}{RGB}{56, 140, 70}
\definecolor{myb}{RGB}{45, 111, 177}
\definecolor{myr}{RGB}{199, 68, 64}
\definecolor{mytheorembg}{HTML}{F2F2F9}
\definecolor{mytheoremfr}{HTML}{00007B}
\definecolor{mylenmabg}{HTML}{FFFAF8}
\definecolor{mylenmafr}{HTML}{983b0f}
\definecolor{mypropbg}{HTML}{f2fbfc}
\definecolor{mypropfr}{HTML}{191971}
\definecolor{myexamplebg}{HTML}{F2FBF8}
\definecolor{myexamplefr}{HTML}{88D6D1}
\definecolor{myexampleti}{HTML}{2A7F7F}
\definecolor{mydefinitbg}{HTML}{E5E5FF}
\definecolor{mydefinitfr}{HTML}{3F3FA3}
\definecolor{notesgreen}{RGB}{0,162,0}
\definecolor{myp}{RGB}{197, 92, 212}
\definecolor{mygr}{HTML}{2C3338}
\definecolor{myred}{RGB}{127,0,0}
\definecolor{myyellow}{RGB}{169,121,69}
\definecolor{myexercisebg}{HTML}{F2FBF8}
\definecolor{myexercisefg}{HTML}{88D6D1}


%%%%%%%%%%%%%%%%%%%%%%%%%%%%
% TCOLORBOX SETUPS
%%%%%%%%%%%%%%%%%%%%%%%%%%%%

\setlength{\parindent}{1cm}
%================================
% THEOREM BOX
%================================

\tcbuselibrary{theorems,skins,hooks}
\newtcbtheorem[number within=section]{Theorem}{Theorem}
{%
	enhanced,
	breakable,
	colback = mytheorembg,
	frame hidden,
	boxrule = 0sp,
	borderline west = {2pt}{0pt}{mytheoremfr},
	sharp corners,
	detach title,
	before upper = \tcbtitle\par\smallskip,
	coltitle = mytheoremfr,
	fonttitle = \bfseries\sffamily,
	description font = \mdseries,
	separator sign none,
	segmentation style={solid, mytheoremfr},
}
{th}

\tcbuselibrary{theorems,skins,hooks}
\newtcbtheorem[number within=chapter]{theorem}{Theorem}
{%
	enhanced,
	breakable,
	colback = mytheorembg,
	frame hidden,
	boxrule = 0sp,
	borderline west = {2pt}{0pt}{mytheoremfr},
	sharp corners,
	detach title,
	before upper = \tcbtitle\par\smallskip,
	coltitle = mytheoremfr,
	fonttitle = \bfseries\sffamily,
	description font = \mdseries,
	separator sign none,
	segmentation style={solid, mytheoremfr},
}
{th}


\tcbuselibrary{theorems,skins,hooks}
\newtcolorbox{Theoremcon}
{%
	enhanced
	,breakable
	,colback = mytheorembg
	,frame hidden
	,boxrule = 0sp
	,borderline west = {2pt}{0pt}{mytheoremfr}
	,sharp corners
	,description font = \mdseries
	,separator sign none
}

%================================
% Corollery
%================================
\tcbuselibrary{theorems,skins,hooks}
\newtcbtheorem[number within=section]{Corollary}{Corollary}
{%
	enhanced
	,breakable
	,colback = myp!10
	,frame hidden
	,boxrule = 0sp
	,borderline west = {2pt}{0pt}{myp!85!black}
	,sharp corners
	,detach title
	,before upper = \tcbtitle\par\smallskip
	,coltitle = myp!85!black
	,fonttitle = \bfseries\sffamily
	,description font = \mdseries
	,separator sign none
	,segmentation style={solid, myp!85!black}
}
{th}
\tcbuselibrary{theorems,skins,hooks}
\newtcbtheorem[number within=chapter]{corollary}{Corollary}
{%
	enhanced
	,breakable
	,colback = myp!10
	,frame hidden
	,boxrule = 0sp
	,borderline west = {2pt}{0pt}{myp!85!black}
	,sharp corners
	,detach title
	,before upper = \tcbtitle\par\smallskip
	,coltitle = myp!85!black
	,fonttitle = \bfseries\sffamily
	,description font = \mdseries
	,separator sign none
	,segmentation style={solid, myp!85!black}
}
{th}


%================================
% LENMA
%================================

\tcbuselibrary{theorems,skins,hooks}
\newtcbtheorem[number within=section]{Lenma}{Lenma}
{%
	enhanced,
	breakable,
	colback = mylenmabg,
	frame hidden,
	boxrule = 0sp,
	borderline west = {2pt}{0pt}{mylenmafr},
	sharp corners,
	detach title,
	before upper = \tcbtitle\par\smallskip,
	coltitle = mylenmafr,
	fonttitle = \bfseries\sffamily,
	description font = \mdseries,
	separator sign none,
	segmentation style={solid, mylenmafr},
}
{th}

\tcbuselibrary{theorems,skins,hooks}
\newtcbtheorem[number within=chapter]{lenma}{Lenma}
{%
	enhanced,
	breakable,
	colback = mylenmabg,
	frame hidden,
	boxrule = 0sp,
	borderline west = {2pt}{0pt}{mylenmafr},
	sharp corners,
	detach title,
	before upper = \tcbtitle\par\smallskip,
	coltitle = mylenmafr,
	fonttitle = \bfseries\sffamily,
	description font = \mdseries,
	separator sign none,
	segmentation style={solid, mylenmafr},
}
{th}


%================================
% PROPOSITION
%================================

\tcbuselibrary{theorems,skins,hooks}
\newtcbtheorem[number within=section]{Prop}{Proposition}
{%
	enhanced,
	breakable,
	colback = mypropbg,
	frame hidden,
	boxrule = 0sp,
	borderline west = {2pt}{0pt}{mypropfr},
	sharp corners,
	detach title,
	before upper = \tcbtitle\par\smallskip,
	coltitle = mypropfr,
	fonttitle = \bfseries\sffamily,
	description font = \mdseries,
	separator sign none,
	segmentation style={solid, mypropfr},
}
{th}

\tcbuselibrary{theorems,skins,hooks}
\newtcbtheorem[number within=chapter]{prop}{Proposition}
{%
	enhanced,
	breakable,
	colback = mypropbg,
	frame hidden,
	boxrule = 0sp,
	borderline west = {2pt}{0pt}{mypropfr},
	sharp corners,
	detach title,
	before upper = \tcbtitle\par\smallskip,
	coltitle = mypropfr,
	fonttitle = \bfseries\sffamily,
	description font = \mdseries,
	separator sign none,
	segmentation style={solid, mypropfr},
}
{th}


%================================
% CLAIM
%================================

\tcbuselibrary{theorems,skins,hooks}
\newtcbtheorem[number within=section]{claim}{Claim}
{%
	enhanced
	,breakable
	,colback = myg!10
	,frame hidden
	,boxrule = 0sp
	,borderline west = {2pt}{0pt}{myg}
	,sharp corners
	,detach title
	,before upper = \tcbtitle\par\smallskip
	,coltitle = myg!85!black
	,fonttitle = \bfseries\sffamily
	,description font = \mdseries
	,separator sign none
	,segmentation style={solid, myg!85!black}
}
{th}



%================================
% Exercise
%================================

\tcbuselibrary{theorems,skins,hooks}
\newtcbtheorem[number within=section]{Exercise}{Exercise}
{%
	enhanced,
	breakable,
	colback = myexercisebg,
	frame hidden,
	boxrule = 0sp,
	borderline west = {2pt}{0pt}{myexercisefg},
	sharp corners,
	detach title,
	before upper = \tcbtitle\par\smallskip,
	coltitle = myexercisefg,
	fonttitle = \bfseries\sffamily,
	description font = \mdseries,
	separator sign none,
	segmentation style={solid, myexercisefg},
}
{th}

\tcbuselibrary{theorems,skins,hooks}
\newtcbtheorem[number within=chapter]{exercise}{Exercise}
{%
	enhanced,
	breakable,
	colback = myexercisebg,
	frame hidden,
	boxrule = 0sp,
	borderline west = {2pt}{0pt}{myexercisefg},
	sharp corners,
	detach title,
	before upper = \tcbtitle\par\smallskip,
	coltitle = myexercisefg,
	fonttitle = \bfseries\sffamily,
	description font = \mdseries,
	separator sign none,
	segmentation style={solid, myexercisefg},
}
{th}

%================================
% EXAMPLE BOX
%================================

\newtcbtheorem[number within=section]{Example}{Example}
{%
	colback = myexamplebg
	,breakable
	,colframe = myexamplefr
	,coltitle = myexampleti
	,boxrule = 1pt
	,sharp corners
	,detach title
	,before upper=\tcbtitle\par\smallskip
	,fonttitle = \bfseries
	,description font = \mdseries
	,separator sign none
	,description delimiters parenthesis
}
{ex}

\newtcbtheorem[number within=chapter]{example}{Example}
{%
	colback = myexamplebg
	,breakable
	,colframe = myexamplefr
	,coltitle = myexampleti
	,boxrule = 1pt
	,sharp corners
	,detach title
	,before upper=\tcbtitle\par\smallskip
	,fonttitle = \bfseries
	,description font = \mdseries
	,separator sign none
	,description delimiters parenthesis
}
{ex}

%================================
% DEFINITION BOX
%================================

\newtcbtheorem[number within=section]{Definition}{Definition}{enhanced,
	before skip=2mm,after skip=2mm, colback=red!5,colframe=red!80!black,boxrule=0.5mm,
	attach boxed title to top left={xshift=1cm,yshift*=1mm-\tcboxedtitleheight}, varwidth boxed title*=-3cm,
	boxed title style={frame code={
					\path[fill=tcbcolback]
					([yshift=-1mm,xshift=-1mm]frame.north west)
					arc[start angle=0,end angle=180,radius=1mm]
					([yshift=-1mm,xshift=1mm]frame.north east)
					arc[start angle=180,end angle=0,radius=1mm];
					\path[left color=tcbcolback!60!black,right color=tcbcolback!60!black,
						middle color=tcbcolback!80!black]
					([xshift=-2mm]frame.north west) -- ([xshift=2mm]frame.north east)
					[rounded corners=1mm]-- ([xshift=1mm,yshift=-1mm]frame.north east)
					-- (frame.south east) -- (frame.south west)
					-- ([xshift=-1mm,yshift=-1mm]frame.north west)
					[sharp corners]-- cycle;
				},interior engine=empty,
		},
	fonttitle=\bfseries,
	title={#2},#1}{def}
\newtcbtheorem[number within=chapter]{definition}{Definition}{enhanced,
	before skip=2mm,after skip=2mm, colback=red!5,colframe=red!80!black,boxrule=0.5mm,
	attach boxed title to top left={xshift=1cm,yshift*=1mm-\tcboxedtitleheight}, varwidth boxed title*=-3cm,
	boxed title style={frame code={
					\path[fill=tcbcolback]
					([yshift=-1mm,xshift=-1mm]frame.north west)
					arc[start angle=0,end angle=180,radius=1mm]
					([yshift=-1mm,xshift=1mm]frame.north east)
					arc[start angle=180,end angle=0,radius=1mm];
					\path[left color=tcbcolback!60!black,right color=tcbcolback!60!black,
						middle color=tcbcolback!80!black]
					([xshift=-2mm]frame.north west) -- ([xshift=2mm]frame.north east)
					[rounded corners=1mm]-- ([xshift=1mm,yshift=-1mm]frame.north east)
					-- (frame.south east) -- (frame.south west)
					-- ([xshift=-1mm,yshift=-1mm]frame.north west)
					[sharp corners]-- cycle;
				},interior engine=empty,
		},
	fonttitle=\bfseries,
	title={#2},#1}{def}



%================================
% Solution BOX
%================================

\makeatletter
\newtcbtheorem{question}{Pregunta}{enhanced,
	breakable,
	colback=white,
	colframe=myb!80!black,
	attach boxed title to top left={yshift*=-\tcboxedtitleheight},
	fonttitle=\bfseries,
	title={#2},
	boxed title size=title,
	boxed title style={%
			sharp corners,
			rounded corners=northwest,
			colback=tcbcolframe,
			boxrule=0pt,
		},
	underlay boxed title={%
			\path[fill=tcbcolframe] (title.south west)--(title.south east)
			to[out=0, in=180] ([xshift=5mm]title.east)--
			(title.center-|frame.east)
			[rounded corners=\kvtcb@arc] |-
			(frame.north) -| cycle;
		},
	#1
}{def}
\makeatother

%================================
% SOLUTION BOX
%================================

\makeatletter
\newtcolorbox{solution}{enhanced,
	breakable,
	colback=white,
	colframe=myg!80!black,
	attach boxed title to top left={yshift*=-\tcboxedtitleheight},
	title=Solution,
	boxed title size=title,
	boxed title style={%
			sharp corners,
			rounded corners=northwest,
			colback=tcbcolframe,
			boxrule=0pt,
		},
	underlay boxed title={%
			\path[fill=tcbcolframe] (title.south west)--(title.south east)
			to[out=0, in=180] ([xshift=5mm]title.east)--
			(title.center-|frame.east)
			[rounded corners=\kvtcb@arc] |-
			(frame.north) -| cycle;
		},
}
\makeatother

%================================
% Question BOX
%================================

\makeatletter
\newtcbtheorem{qstion}{Pregunta}{enhanced,
	breakable,
	colback=white,
	colframe=mygr,
	attach boxed title to top left={yshift*=-\tcboxedtitleheight},
	fonttitle=\bfseries,
	title={#2},
	boxed title size=title,
	boxed title style={%
			sharp corners,
			rounded corners=northwest,
			colback=tcbcolframe,
			boxrule=0pt,
		},
	underlay boxed title={%
			\path[fill=tcbcolframe] (title.south west)--(title.south east)
			to[out=0, in=180] ([xshift=5mm]title.east)--
			(title.center-|frame.east)
			[rounded corners=\kvtcb@arc] |-
			(frame.north) -| cycle;
		},
	#1
}{def}
\makeatother

\newtcbtheorem[number within=chapter]{wconc}{Wrong Concept}{
	breakable,
	enhanced,
	colback=white,
	colframe=myr,
	arc=0pt,
	outer arc=0pt,
	fonttitle=\bfseries\sffamily\large,
	colbacktitle=myr,
	attach boxed title to top left={},
	boxed title style={
			enhanced,
			skin=enhancedfirst jigsaw,
			arc=3pt,
			bottom=0pt,
			interior style={fill=myr}
		},
	#1
}{def}



%================================
% NOTE BOX
%================================

\usetikzlibrary{arrows,calc,shadows.blur}
\tcbuselibrary{skins}
\newtcolorbox{note}[1][]{%
	enhanced jigsaw,
	colback=gray!20!white,%
	colframe=gray!80!black,
	size=small,
	boxrule=1pt,
	title=\textbf{Nota:-},
	halign title=flush center,
	coltitle=black,
	breakable,
	drop shadow=black!50!white,
	attach boxed title to top left={xshift=1cm,yshift=-\tcboxedtitleheight/2,yshifttext=-\tcboxedtitleheight/2},
	minipage boxed title=1.5cm,
	boxed title style={%
			colback=white,
			size=fbox,
			boxrule=1pt,
			boxsep=2pt,
			underlay={%
					\coordinate (dotA) at ($(interior.west) + (-0.5pt,0)$);
					\coordinate (dotB) at ($(interior.east) + (0.5pt,0)$);
					\begin{scope}
						\clip (interior.north west) rectangle ([xshift=3ex]interior.east);
						\filldraw [white, blur shadow={shadow opacity=60, shadow yshift=-.75ex}, rounded corners=2pt] (interior.north west) rectangle (interior.south east);
					\end{scope}
					\begin{scope}[gray!80!black]
						\fill (dotA) circle (2pt);
						\fill (dotB) circle (2pt);
					\end{scope}
				},
		},
	#1,
}

%%%%%%%%%%%%%%%%%%%%%%%%%%%%%%
% SELF MADE COMMANDS
%%%%%%%%%%%%%%%%%%%%%%%%%%%%%%


\newcommand{\thm}[2]{\begin{Theorem}{#1}{}#2\end{Theorem}}
\newcommand{\cor}[2]{\begin{Corollary}{#1}{}#2\end{Corollary}}
\newcommand{\mlenma}[2]{\begin{Lenma}{#1}{}#2\end{Lenma}}
\newcommand{\mprop}[2]{\begin{Prop}{#1}{}#2\end{Prop}}
\newcommand{\clm}[3]{\begin{claim}{#1}{#2}#3\end{claim}}
\newcommand{\wc}[2]{\begin{wconc}{#1}{}\setlength{\parindent}{1cm}#2\end{wconc}}
\newcommand{\thmcon}[1]{\begin{Theoremcon}{#1}\end{Theoremcon}}
\newcommand{\ex}[2]{\begin{Example}{#1}{}#2\end{Example}}
\newcommand{\dfn}[2]{\begin{Definition}[colbacktitle=red!75!black]{#1}{}#2\end{Definition}}
\newcommand{\dfnc}[2]{\begin{definition}[colbacktitle=red!75!black]{#1}{}#2\end{definition}}
\newcommand{\qs}[2]{\begin{question}{#1}{}#2\end{question}}
\newcommand{\pf}[2]{\begin{myproof}[#1]#2\end{myproof}}
\NewDocumentCommand{\nt}{+m}{\begin{note}#1\end{note}}

\newcommand*\circled[1]{\tikz[baseline=(char.base)]{
		\node[shape=circle,draw,inner sep=1pt] (char) {#1};}}
\newcommand\getcurrentref[1]{%
	\ifnumequal{\value{#1}}{0}
	{??}
	{\the\value{#1}}%
}
\newcommand{\getCurrentSectionNumber}{\getcurrentref{section}}
\newenvironment{myproof}[1][\proofname]{%
	\proof[\bfseries #1: ]%
}{\endproof}

\newcommand{\mclm}[2]{\begin{myclaim}[#1]#2\end{myclaim}}
\newenvironment{myclaim}[1][\claimname]{\proof[\bfseries #1: ]}{}

\newcounter{mylabelcounter}

\makeatletter
\newcommand{\setword}[2]{%
	\phantomsection
	#1\def\@currentlabel{\unexpanded{#1}}\label{#2}%
}
\makeatother




\tikzset{
	symbol/.style={
			draw=none,
			every to/.append style={
					edge node={node [sloped, allow upside down, auto=false]{$#1$}}}
		}
}


% deliminators
\DeclarePairedDelimiter{\abs}{\lvert}{\rvert}
\DeclarePairedDelimiter{\norm}{\lVert}{\rVert}

\DeclarePairedDelimiter{\ceil}{\lceil}{\rceil}
\DeclarePairedDelimiter{\floor}{\lfloor}{\rfloor}
\DeclarePairedDelimiter{\round}{\lfloor}{\rceil}

\newsavebox\diffdbox
\newcommand{\slantedromand}{{\mathpalette\makesl{d}}}
\newcommand{\makesl}[2]{%
\begingroup
\sbox{\diffdbox}{$\mathsurround=0pt#1\mathrm{#2}$}%
\pdfsave
\pdfsetmatrix{1 0 0.2 1}%
\rlap{\usebox{\diffdbox}}%
\pdfrestore
\hskip\wd\diffdbox
\endgroup
}
\newcommand{\dd}[1][]{\ensuremath{\mathop{}\!\ifstrempty{#1}{%
\slantedromand\@ifnextchar^{\hspace{0.2ex}}{\hspace{0.1ex}}}%
{\slantedromand\hspace{0.2ex}^{#1}}}}
\ProvideDocumentCommand\dv{o m g}{%
  \ensuremath{%
    \IfValueTF{#3}{%
      \IfNoValueTF{#1}{%
        \frac{\dd #2}{\dd #3}%
      }{%
        \frac{\dd^{#1} #2}{\dd #3^{#1}}%
      }%
    }{%
      \IfNoValueTF{#1}{%
        \frac{\dd}{\dd #2}%
      }{%
        \frac{\dd^{#1}}{\dd #2^{#1}}%
      }%
    }%
  }%
}
\providecommand*{\pdv}[3][]{\frac{\partial^{#1}#2}{\partial#3^{#1}}}
%  - others
\DeclareMathOperator{\Lap}{\mathcal{L}}
\DeclareMathOperator{\Var}{Var} % varience
\DeclareMathOperator{\Cov}{Cov} % covarience
\DeclareMathOperator{\E}{E} % expected

% Since the amsthm package isn't loaded

% I prefer the slanted \leq
\let\oldleq\leq % save them in case they're every wanted
\let\oldgeq\geq
\renewcommand{\leq}{\leqslant}
\renewcommand{\geq}{\geqslant}

% % redefine matrix env to allow for alignment, use r as default
% \renewcommand*\env@matrix[1][r]{\hskip -\arraycolsep
%     \let\@ifnextchar\new@ifnextchar
%     \array{*\c@MaxMatrixCols #1}}


%\usepackage{framed}
%\usepackage{titletoc}
%\usepackage{etoolbox}
%\usepackage{lmodern}


%\patchcmd{\tableofcontents}{\contentsname}{\sffamily\contentsname}{}{}

%\renewenvironment{leftbar}
%{\def\FrameCommand{\hspace{6em}%
%		{\color{myyellow}\vrule width 2pt depth 6pt}\hspace{1em}}%
%	\MakeFramed{\parshape 1 0cm \dimexpr\textwidth-6em\relax\FrameRestore}\vskip2pt%
%}
%{\endMakeFramed}

%\titlecontents{chapter}
%[0em]{\vspace*{2\baselineskip}}
%{\parbox{4.5em}{%
%		\hfill\Huge\sffamily\bfseries\color{myred}\thecontentspage}%
%	\vspace*{-2.3\baselineskip}\leftbar\textsc{\small\chaptername~\thecontentslabel}\\\sffamily}
%{}{\endleftbar}
%\titlecontents{section}
%[8.4em]
%{\sffamily\contentslabel{3em}}{}{}
%{\hspace{0.5em}\nobreak\itshape\color{myred}\contentspage}
%\titlecontents{subsection}
%[8.4em]
%{\sffamily\contentslabel{3em}}{}{}  
%{\hspace{0.5em}\nobreak\itshape\color{myred}\contentspage}



%%%%%%%%%%%%%%%%%%%%%%%%%%%%%%%%%%%%%%%%%%%
% TABLE OF CONTENTS
%%%%%%%%%%%%%%%%%%%%%%%%%%%%%%%%%%%%%%%%%%%

\usepackage{tikz}
\definecolor{doc}{RGB}{0,60,110}
\usepackage{titletoc}
\contentsmargin{0cm}
\titlecontents{chapter}[3.7pc]
{\addvspace{30pt}%
	\begin{tikzpicture}[remember picture, overlay]%
		\draw[fill=doc!60,draw=doc!60] (-7,-.1) rectangle (-0.9,.5);%
		\pgftext[left,x=-3.5cm,y=0.2cm]{\color{white}\Large\sc\bfseries Módulo\ \thecontentslabel};%
	\end{tikzpicture}\color{doc!60}\large\sc\bfseries}%
{}
{}
{\;\titlerule\;\large\sc\bfseries Page \thecontentspage
	\begin{tikzpicture}[remember picture, overlay]
		\draw[fill=doc!60,draw=doc!60] (2pt,0) rectangle (4,0.1pt);
	\end{tikzpicture}}%
\titlecontents{section}[3.7pc]
{\addvspace{2pt}}
{\contentslabel[\thecontentslabel]{2pc}}
{}
{\hfill\small \thecontentspage}
[]

\titlecontents*{subsection}[4.7pc] % Mayor indentación para subsecciones
{\addvspace{1pt}\small} % Espaciado más limpio
{\thecontentslabel\quad} % Mostrar número de subsección (1.1.1, 1.1.2)
{}
{\hfill\small\thecontentspage\\} % Alinear el número de página correctamente


\makeatletter
\renewcommand{\tableofcontents}{%
	\chapter*{%
	  \vspace*{-20\p@}%
	  \begin{tikzpicture}[remember picture, overlay]%
		  \pgftext[right,x=15cm,y=0.2cm]{\color{doc!60}\Huge\sc\bfseries \contentsname};%
		  \draw[fill=doc!60,draw=doc!60] (13.5,-.75) rectangle (20,1);%
		  \clip (13.4,-.75) rectangle (20,1);
		  \pgftext[right,x=15cm,y=0.2cm]{\color{white}\Huge\sc\bfseries \contentsname};%
	  \end{tikzpicture}}%
	\@starttoc{toc}}
\makeatother

 % Configuración y paquetes de diseño
%From M275 "Topology" at SJSU
\newcommand{\id}{\mathrm{id}}
\newcommand{\taking}[1]{\xrightarrow{#1}}
\newcommand{\inv}{^{-1}}

%From M170 "Introduction to Graph Theory" at SJSU
\DeclareMathOperator{\diam}{diam}
\DeclareMathOperator{\ord}{ord}
\newcommand{\defeq}{\overset{\mathrm{def}}{=}}

%From the USAMO .tex files
\newcommand{\ts}{\textsuperscript}
\newcommand{\dg}{^\circ}
\newcommand{\ii}{\item}

% % From Math 55 and Math 145 at Harvard
% \newenvironment{subproof}[1][Proof]{%
% \begin{proof}[#1] \renewcommand{\qedsymbol}{$\blacksquare$}}%
% {\end{proof}}

\newcommand{\liff}{\leftrightarrow}
\newcommand{\lthen}{\rightarrow}
\newcommand{\opname}{\operatorname}
\newcommand{\surjto}{\twoheadrightarrow}
\newcommand{\injto}{\hookrightarrow}
\newcommand{\On}{\mathrm{On}} % ordinals
\DeclareMathOperator{\img}{im} % Image
\DeclareMathOperator{\Img}{Im} % Image
\DeclareMathOperator{\coker}{coker} % Cokernel
\DeclareMathOperator{\Coker}{Coker} % Cokernel
\DeclareMathOperator{\Ker}{Ker} % Kernel
\DeclareMathOperator{\rank}{rank}
\DeclareMathOperator{\Spec}{Spec} % spectrum
\DeclareMathOperator{\Tr}{Tr} % trace
\DeclareMathOperator{\pr}{pr} % projection
\DeclareMathOperator{\ext}{ext} % extension
\DeclareMathOperator{\pred}{pred} % predecessor
\DeclareMathOperator{\dom}{dom} % domain
\DeclareMathOperator{\ran}{ran} % range
\DeclareMathOperator{\Hom}{Hom} % homomorphism
\DeclareMathOperator{\Mor}{Mor} % morphisms
\DeclareMathOperator{\End}{End} % endomorphism

\newcommand{\eps}{\epsilon}
\newcommand{\veps}{\varepsilon}
\newcommand{\ol}{\overline}
\newcommand{\ul}{\underline}
\newcommand{\wt}{\widetilde}
\newcommand{\wh}{\widehat}
\newcommand{\vocab}[1]{\textbf{\color{blue} #1}}
\providecommand{\half}{\frac{1}{2}}
\newcommand{\dang}{\measuredangle} %% Directed angle
\newcommand{\ray}[1]{\overrightarrow{#1}}
\newcommand{\seg}[1]{\overline{#1}}
\newcommand{\arc}[1]{\wideparen{#1}}
\DeclareMathOperator{\cis}{cis}
\DeclareMathOperator*{\lcm}{lcm}
\DeclareMathOperator*{\argmin}{arg min}
\DeclareMathOperator*{\argmax}{arg max}
\newcommand{\cycsum}{\sum_{\mathrm{cyc}}}
\newcommand{\symsum}{\sum_{\mathrm{sym}}}
\newcommand{\cycprod}{\prod_{\mathrm{cyc}}}
\newcommand{\symprod}{\prod_{\mathrm{sym}}}
\newcommand{\Qed}{\begin{flushright}\qed\end{flushright}}
\newcommand{\parinn}{\setlength{\parindent}{1cm}}
\newcommand{\parinf}{\setlength{\parindent}{0cm}}
% \newcommand{\norm}{\|\cdot\|}
\newcommand{\inorm}{\norm_{\infty}}
\newcommand{\opensets}{\{V_{\alpha}\}_{\alpha\in I}}
\newcommand{\oset}{V_{\alpha}}
\newcommand{\opset}[1]{V_{\alpha_{#1}}}
\newcommand{\lub}{\text{lub}}
\newcommand{\del}[2]{\frac{\partial #1}{\partial #2}}
\newcommand{\Del}[3]{\frac{\partial^{#1} #2}{\partial^{#1} #3}}
\newcommand{\deld}[2]{\dfrac{\partial #1}{\partial #2}}
\newcommand{\Deld}[3]{\dfrac{\partial^{#1} #2}{\partial^{#1} #3}}
\newcommand{\lm}{\lambda}
\newcommand{\uin}{\mathbin{\rotatebox[origin=c]{90}{$\in$}}}
\newcommand{\usubset}{\mathbin{\rotatebox[origin=c]{90}{$\subset$}}}
\newcommand{\lt}{\left}
\newcommand{\rt}{\right}
\newcommand{\bs}[1]{\boldsymbol{#1}}
\newcommand{\exs}{\exists}
\newcommand{\st}{\strut}
\newcommand{\dps}[1]{\displaystyle{#1}}

\newcommand{\sol}{\setlength{\parindent}{0cm}\textbf{\textit{Solution:}}\setlength{\parindent}{1cm} }
\newcommand{\solve}[1]{\setlength{\parindent}{0cm}\textbf{\textit{Solución: }}\setlength{\parindent}{1cm}#1 \Qed}

% Things Lie
\newcommand{\kb}{\mathfrak b}
\newcommand{\kg}{\mathfrak g}
\newcommand{\kh}{\mathfrak h}
\newcommand{\kn}{\mathfrak n}
\newcommand{\ku}{\mathfrak u}
\newcommand{\kz}{\mathfrak z}
\DeclareMathOperator{\Ext}{Ext} % Ext functor
\DeclareMathOperator{\Tor}{Tor} % Tor functor
\newcommand{\gl}{\opname{\mathfrak{gl}}} % frak gl group
\renewcommand{\sl}{\opname{\mathfrak{sl}}} % frak sl group chktex 6

% More script letters etc.
\newcommand{\SA}{\mathcal A}
\newcommand{\SB}{\mathcal B}
\newcommand{\SC}{\mathcal C}
\newcommand{\SF}{\mathcal F}
\newcommand{\SG}{\mathcal G}
\newcommand{\SH}{\mathcal H}
\newcommand{\OO}{\mathcal O}

\newcommand{\SCA}{\mathscr A}
\newcommand{\SCB}{\mathscr B}
\newcommand{\SCC}{\mathscr C}
\newcommand{\SCD}{\mathscr D}
\newcommand{\SCE}{\mathscr E}
\newcommand{\SCF}{\mathscr F}
\newcommand{\SCG}{\mathscr G}
\newcommand{\SCH}{\mathscr H}

% Mathfrak primes
\newcommand{\km}{\mathfrak m}
\newcommand{\kp}{\mathfrak p}
\newcommand{\kq}{\mathfrak q}

% number sets
\newcommand{\RR}[1][]{\ensuremath{\ifstrempty{#1}{\mathbb{R}}{\mathbb{R}^{#1}}}}
\newcommand{\NN}[1][]{\ensuremath{\ifstrempty{#1}{\mathbb{N}}{\mathbb{N}^{#1}}}}
\newcommand{\ZZ}[1][]{\ensuremath{\ifstrempty{#1}{\mathbb{Z}}{\mathbb{Z}^{#1}}}}
\newcommand{\QQ}[1][]{\ensuremath{\ifstrempty{#1}{\mathbb{Q}}{\mathbb{Q}^{#1}}}}
\newcommand{\CC}[1][]{\ensuremath{\ifstrempty{#1}{\mathbb{C}}{\mathbb{C}^{#1}}}}
\newcommand{\PP}[1][]{\ensuremath{\ifstrempty{#1}{\mathbb{P}}{\mathbb{P}^{#1}}}}
\newcommand{\HH}[1][]{\ensuremath{\ifstrempty{#1}{\mathbb{H}}{\mathbb{H}^{#1}}}}
\newcommand{\FF}[1][]{\ensuremath{\ifstrempty{#1}{\mathbb{F}}{\mathbb{F}^{#1}}}}
% expected value
\newcommand{\EE}{\ensuremath{\mathbb{E}}}
\newcommand{\charin}{\text{ char }}
\DeclareMathOperator{\sign}{sign}
\DeclareMathOperator{\Aut}{Aut}
\DeclareMathOperator{\Inn}{Inn}
\DeclareMathOperator{\Syl}{Syl}
\DeclareMathOperator{\Gal}{Gal}
\DeclareMathOperator{\GL}{GL} % General linear group
\DeclareMathOperator{\SL}{SL} % Special linear group

%---------------------------------------
% BlackBoard Math Fonts :-
%---------------------------------------

%Captital Letters
\newcommand{\bbA}{\mathbb{A}}	\newcommand{\bbB}{\mathbb{B}}
\newcommand{\bbC}{\mathbb{C}}	\newcommand{\bbD}{\mathbb{D}}
\newcommand{\bbE}{\mathbb{E}}	\newcommand{\bbF}{\mathbb{F}}
\newcommand{\bbG}{\mathbb{G}}	\newcommand{\bbH}{\mathbb{H}}
\newcommand{\bbI}{\mathbb{I}}	\newcommand{\bbJ}{\mathbb{J}}
\newcommand{\bbK}{\mathbb{K}}	\newcommand{\bbL}{\mathbb{L}}
\newcommand{\bbM}{\mathbb{M}}	\newcommand{\bbN}{\mathbb{N}}
\newcommand{\bbO}{\mathbb{O}}	\newcommand{\bbP}{\mathbb{P}}
\newcommand{\bbQ}{\mathbb{Q}}	\newcommand{\bbR}{\mathbb{R}}
\newcommand{\bbS}{\mathbb{S}}	\newcommand{\bbT}{\mathbb{T}}
\newcommand{\bbU}{\mathbb{U}}	\newcommand{\bbV}{\mathbb{V}}
\newcommand{\bbW}{\mathbb{W}}	\newcommand{\bbX}{\mathbb{X}}
\newcommand{\bbY}{\mathbb{Y}}	\newcommand{\bbZ}{\mathbb{Z}}

%---------------------------------------
% MathCal Fonts :-
%---------------------------------------

%Captital Letters
\newcommand{\mcA}{\mathcal{A}}	\newcommand{\mcB}{\mathcal{B}}
\newcommand{\mcC}{\mathcal{C}}	\newcommand{\mcD}{\mathcal{D}}
\newcommand{\mcE}{\mathcal{E}}	\newcommand{\mcF}{\mathcal{F}}
\newcommand{\mcG}{\mathcal{G}}	\newcommand{\mcH}{\mathcal{H}}
\newcommand{\mcI}{\mathcal{I}}	\newcommand{\mcJ}{\mathcal{J}}
\newcommand{\mcK}{\mathcal{K}}	\newcommand{\mcL}{\mathcal{L}}
\newcommand{\mcM}{\mathcal{M}}	\newcommand{\mcN}{\mathcal{N}}
\newcommand{\mcO}{\mathcal{O}}	\newcommand{\mcP}{\mathcal{P}}
\newcommand{\mcQ}{\mathcal{Q}}	\newcommand{\mcR}{\mathcal{R}}
\newcommand{\mcS}{\mathcal{S}}	\newcommand{\mcT}{\mathcal{T}}
\newcommand{\mcU}{\mathcal{U}}	\newcommand{\mcV}{\mathcal{V}}
\newcommand{\mcW}{\mathcal{W}}	\newcommand{\mcX}{\mathcal{X}}
\newcommand{\mcY}{\mathcal{Y}}	\newcommand{\mcZ}{\mathcal{Z}}


%---------------------------------------
% Bold Math Fonts :-
%---------------------------------------

%Captital Letters
\newcommand{\bmA}{\boldsymbol{A}}	\newcommand{\bmB}{\boldsymbol{B}}
\newcommand{\bmC}{\boldsymbol{C}}	\newcommand{\bmD}{\boldsymbol{D}}
\newcommand{\bmE}{\boldsymbol{E}}	\newcommand{\bmF}{\boldsymbol{F}}
\newcommand{\bmG}{\boldsymbol{G}}	\newcommand{\bmH}{\boldsymbol{H}}
\newcommand{\bmI}{\boldsymbol{I}}	\newcommand{\bmJ}{\boldsymbol{J}}
\newcommand{\bmK}{\boldsymbol{K}}	\newcommand{\bmL}{\boldsymbol{L}}
\newcommand{\bmM}{\boldsymbol{M}}	\newcommand{\bmN}{\boldsymbol{N}}
\newcommand{\bmO}{\boldsymbol{O}}	\newcommand{\bmP}{\boldsymbol{P}}
\newcommand{\bmQ}{\boldsymbol{Q}}	\newcommand{\bmR}{\boldsymbol{R}}
\newcommand{\bmS}{\boldsymbol{S}}	\newcommand{\bmT}{\boldsymbol{T}}
\newcommand{\bmU}{\boldsymbol{U}}	\newcommand{\bmV}{\boldsymbol{V}}
\newcommand{\bmW}{\boldsymbol{W}}	\newcommand{\bmX}{\boldsymbol{X}}
\newcommand{\bmY}{\boldsymbol{Y}}	\newcommand{\bmZ}{\boldsymbol{Z}}
%Small Letters
\newcommand{\bma}{\boldsymbol{a}}	\newcommand{\bmb}{\boldsymbol{b}}
\newcommand{\bmc}{\boldsymbol{c}}	\newcommand{\bmd}{\boldsymbol{d}}
\newcommand{\bme}{\boldsymbol{e}}	\newcommand{\bmf}{\boldsymbol{f}}
\newcommand{\bmg}{\boldsymbol{g}}	\newcommand{\bmh}{\boldsymbol{h}}
\newcommand{\bmi}{\boldsymbol{i}}	\newcommand{\bmj}{\boldsymbol{j}}
\newcommand{\bmk}{\boldsymbol{k}}	\newcommand{\bml}{\boldsymbol{l}}
\newcommand{\bmm}{\boldsymbol{m}}	\newcommand{\bmn}{\boldsymbol{n}}
\newcommand{\bmo}{\boldsymbol{o}}	\newcommand{\bmp}{\boldsymbol{p}}
\newcommand{\bmq}{\boldsymbol{q}}	\newcommand{\bmr}{\boldsymbol{r}}
\newcommand{\bms}{\boldsymbol{s}}	\newcommand{\bmt}{\boldsymbol{t}}
\newcommand{\bmu}{\boldsymbol{u}}	\newcommand{\bmv}{\boldsymbol{v}}
\newcommand{\bmw}{\boldsymbol{w}}	\newcommand{\bmx}{\boldsymbol{x}}
\newcommand{\bmy}{\boldsymbol{y}}	\newcommand{\bmz}{\boldsymbol{z}}

%---------------------------------------
% Scr Math Fonts :-
%---------------------------------------

\newcommand{\sA}{{\mathscr{A}}}   \newcommand{\sB}{{\mathscr{B}}}
\newcommand{\sC}{{\mathscr{C}}}   \newcommand{\sD}{{\mathscr{D}}}
\newcommand{\sE}{{\mathscr{E}}}   \newcommand{\sF}{{\mathscr{F}}}
\newcommand{\sG}{{\mathscr{G}}}   \newcommand{\sH}{{\mathscr{H}}}
\newcommand{\sI}{{\mathscr{I}}}   \newcommand{\sJ}{{\mathscr{J}}}
\newcommand{\sK}{{\mathscr{K}}}   \newcommand{\sL}{{\mathscr{L}}}
\newcommand{\sM}{{\mathscr{M}}}   \newcommand{\sN}{{\mathscr{N}}}
\newcommand{\sO}{{\mathscr{O}}}   \newcommand{\sP}{{\mathscr{P}}}
\newcommand{\sQ}{{\mathscr{Q}}}   \newcommand{\sR}{{\mathscr{R}}}
\newcommand{\sS}{{\mathscr{S}}}   \newcommand{\sT}{{\mathscr{T}}}
\newcommand{\sU}{{\mathscr{U}}}   \newcommand{\sV}{{\mathscr{V}}}
\newcommand{\sW}{{\mathscr{W}}}   \newcommand{\sX}{{\mathscr{X}}}
\newcommand{\sY}{{\mathscr{Y}}}   \newcommand{\sZ}{{\mathscr{Z}}}


%---------------------------------------
% Math Fraktur Font
%---------------------------------------

%Captital Letters
\newcommand{\mfA}{\mathfrak{A}}	\newcommand{\mfB}{\mathfrak{B}}
\newcommand{\mfC}{\mathfrak{C}}	\newcommand{\mfD}{\mathfrak{D}}
\newcommand{\mfE}{\mathfrak{E}}	\newcommand{\mfF}{\mathfrak{F}}
\newcommand{\mfG}{\mathfrak{G}}	\newcommand{\mfH}{\mathfrak{H}}
\newcommand{\mfI}{\mathfrak{I}}	\newcommand{\mfJ}{\mathfrak{J}}
\newcommand{\mfK}{\mathfrak{K}}	\newcommand{\mfL}{\mathfrak{L}}
\newcommand{\mfM}{\mathfrak{M}}	\newcommand{\mfN}{\mathfrak{N}}
\newcommand{\mfO}{\mathfrak{O}}	\newcommand{\mfP}{\mathfrak{P}}
\newcommand{\mfQ}{\mathfrak{Q}}	\newcommand{\mfR}{\mathfrak{R}}
\newcommand{\mfS}{\mathfrak{S}}	\newcommand{\mfT}{\mathfrak{T}}
\newcommand{\mfU}{\mathfrak{U}}	\newcommand{\mfV}{\mathfrak{V}}
\newcommand{\mfW}{\mathfrak{W}}	\newcommand{\mfX}{\mathfrak{X}}
\newcommand{\mfY}{\mathfrak{Y}}	\newcommand{\mfZ}{\mathfrak{Z}}
%Small Letters
\newcommand{\mfa}{\mathfrak{a}}	\newcommand{\mfb}{\mathfrak{b}}
\newcommand{\mfc}{\mathfrak{c}}	\newcommand{\mfd}{\mathfrak{d}}
\newcommand{\mfe}{\mathfrak{e}}	\newcommand{\mff}{\mathfrak{f}}
\newcommand{\mfg}{\mathfrak{g}}	\newcommand{\mfh}{\mathfrak{h}}
\newcommand{\mfi}{\mathfrak{i}}	\newcommand{\mfj}{\mathfrak{j}}
\newcommand{\mfk}{\mathfrak{k}}	\newcommand{\mfl}{\mathfrak{l}}
\newcommand{\mfm}{\mathfrak{m}}	\newcommand{\mfn}{\mathfrak{n}}
\newcommand{\mfo}{\mathfrak{o}}	\newcommand{\mfp}{\mathfrak{p}}
\newcommand{\mfq}{\mathfrak{q}}	\newcommand{\mfr}{\mathfrak{r}}
\newcommand{\mfs}{\mathfrak{s}}	\newcommand{\mft}{\mathfrak{t}}
\newcommand{\mfu}{\mathfrak{u}}	\newcommand{\mfv}{\mathfrak{v}}
\newcommand{\mfw}{\mathfrak{w}}	\newcommand{\mfx}{\mathfrak{x}}
\newcommand{\mfy}{\mathfrak{y}}	\newcommand{\mfz}{\mathfrak{z}}


\usepackage{booktabs}
\usepackage{siunitx}
\usepackage{adjustbox}
\usepackage{pdflscape}
\usepackage[backend=biber, style=apa]{biblatex}
\usepackage{caption}
\usepackage{multirow}
\usepackage{csquotes}
\usepackage{subcaption}
\usepackage{float}
\usetikzlibrary{positioning}

\addto\captionsspanish{
  \renewcommand{\contentsname}{Índice General}
}

\graphicspath{ {./Doc/Fotos/} {./} }
\addbibresource{Referencias.bib}

\definecolor{darkgreen}{rgb}{0,0.6,0}

\title{
    \includegraphics[width=0.35\textwidth]{escudo.png}\\[1.5cm]
    {\Huge\textbf{PROYECTO INDUSTRIAL}}\\[0.8em]
    {\LARGE Biorrefinería de Hueso de Aceituna y Poda de Olivo}\\[0.8em]
    {\Large\textbf{Documento nº 1: Estudio de Mercado, Capacidad y Localización}}\\[1.5em]
}
\author{\Large\textbf{Christian Pérez Hita}\\[0.5em] \large Grado en Ingeniería Química / Máster en Ingeniería Química \\ \large Universidad de Granada (UGR)}
\date{\today}

\begin{document}
\spanishdeactivate{.<>}
\setcounter{tocdepth}{3}
\maketitle

\newpage
\pdfbookmark[section]{\contentsname}{contenidos}
\tableofcontents
\pagebreak

%=============================================================================
\chapter{Introducción y Marco del Proyecto}

El sector oleícola en España genera anualmente entre \num{600000} y \num{1000000} de toneladas de hueso de aceituna según la campaña agronómica, representando aproximadamente un \SI{10}{\percent} en masa del fruto procesado \autocite{MAPA2025, AICA2026}. Actualmente, cerca de un \SI{80}{\percent} de este recurso biomásico se comercializa para aprovechamiento térmico convencional en calderas residenciales e industriales, mientras que el \SI{20}{\percent} restante se autoconsume en las propias almazaras \autocite{AVEBIOM2026}. El precio del hueso a granel se sitúa en el mercado térmico en torno a 246--259~€/t con transporte e IVA incluidos, marcando el coste de oportunidad real de la materia prima.

El presente proyecto plantea la implantación de una \textbf{biorrefinería multiproducto de alta eficiencia} capaz de valorizar integralmente la biomasa residual del olivar: el hueso de aceituna, la hoja de olivo y la madera de poda procedente del corte bianual del olivar \autocite{Jafari-Gohargan2026}. Mediante un esquema de fraccionamiento secuencial (hidrólisis ácida suave y organosolv), se sintetiza bioetanol de segunda generación (2G), lignina técnica pura, dióxido de carbono biogénico de alta pureza, biogás de autoconsumo y biofertilizante \autocite{Padilla-Rascon2022, ElFessikh2025}.

\section{Fuentes de Información y Bases de Datos Sectoriales}
Para garantizar el rigor metodológico, los parámetros de mercado, costes y balances se han triangulado a través de las siguientes fuentes oficiales:
\begin{itemize}
    \item \textbf{SABI (Informa / Bureau van Dijk) \autocite{SABI2026}:} Geolocalización y balance económico de almazaras mediante el código \textbf{CNAE 1043} (Fabricación de aceite de oliva), entamadoras con \textbf{CNAE 1039} y productores químicos con \textbf{CNAE 2014}. Permite cartografiar la generación de hueso y auditar la solvencia financiera de proveedores y competidores.
    \item \textbf{AICA y MAPA \autocite{AICA2026, MAPA2025}:} Datos oficiales del Anuario de Estadística y la encuesta ESYRCE, aportando aforos mensuales de entradas de aceituna, producción de aceite y aceituna de mesa por provincias y comarcas.
    \item \textbf{SIAE (Junta de Andalucía) \autocite{SIAE2025}:} Seguimiento del mercado oleícola andaluz, donde se concentran más de \SI{30000}{\tonne/a\text{\~{n}}o} de hueso procedente en exclusiva del aderezo y entamado de aceituna de mesa en el Valle del Guadalquivir y Poniente.
    \item \textbf{AVEBIOM \autocite{AVEBIOM2026}:} Índices de precios de biocombustibles sólidos y hueso de aceituna para evaluar la presión competitiva del canal térmico.
    \item \textbf{CNMC y CORES \autocite{CORES2026}:} Informes estadísticos sobre biocarburantes, certificaciones de venta de bioetanol y cuotas de mezcla obligatorias exigidas a los agentes petrolíferos (*off-takers* como Repsol, Moeve/Cepsa y BP).
    \item \textbf{Eurostat e INE \autocite{Eurostat2025, INE2025}:} Balances comunitarios de bioetanol (códigos \texttt{nrg\_cb\_rw} y \texttt{apro\_cpsh1}) y estadísticas de costes salariales y fletes terrestres industriales.
\end{itemize}

\section{Composición Fisicoquímica de la Materia Prima}
El hueso de aceituna se caracteriza por una elevada densidad aparente y un contenido en humedad relativamente reducido (\SI{8}{\percent}), con una granulometría habitual entre \SI{1}{\milli\meter} y \SI{3.35}{\milli\meter}. Su caracterización en base seca (b.s.) se recoge en la Tabla~\ref{tab:composicion_hueso} \autocite{Padilla-Rascon2022, MacArez2025}.

\begin{table}[H]
    \centering
    \caption{Composición química del hueso de aceituna en base seca.}
    \adjustbox{max width=0.85\textwidth}{%
        \renewcommand{\arraystretch}{1.25}
        \begin{tabular}{lc p{7.5cm}}
            \toprule
            \textbf{Componente} & \textbf{Fracción másica (\% b.s.)} & \textbf{Destino en el Proceso de Biorrefinería}                                                                 \\
            \midrule
            Lignina total       & \SI{35.6}{\percent}                & Fraccionamiento Organosolv $\rightarrow$ Lignina técnica de alta pureza y combustible de caldera                \\
            Hemicelulosa        & \SI{26.0}{\percent}                & Hidrólisis ácida suave $\rightarrow$ Licor rico en pentosas (xilosa $>$ \SI{80}{\percent}) para fermentación 2G \\
            Celulosa            & \SI{20.9}{\percent}                & Sacarificación enzimática + SSF $\rightarrow$ Hexosas (glucosa) para fermentación a bioetanol                   \\
            Extractivos         & \SI{6.3}{\percent}                 & Recuperación de compuestos bioactivos / valorización energética                                                 \\
            Grupos acetilo      & \SI{5.9}{\percent}                 & Hidrólisis a ácido acético (recuperable o biodigerible)                                                         \\
            Cenizas             & \SI{0.6}{\percent}                 & Purgas de combustión / formulación mineral                                                                      \\
            \bottomrule
        \end{tabular}%
    }
    \label{tab:composicion_hueso}
\end{table}

%=============================================================================
\chapter{Estudio de Escala y Capacidad de la Planta}

\section{Escenarios de Capacidad y Selección de Escala}
A partir del modelo tecnoeconómico de referencia de \textcite{Jafari-Gohargan2026}, se contrastan los distintos escenarios de capacidad productiva de biomasa de olivar (Tabla~\ref{tab:escenarios_referencia}).

\begin{table}[H]
    \centering
    \caption{Escenarios de capacidad de tratamiento en base seca.}
    \adjustbox{max width=0.95\textwidth}{%
        \renewcommand{\arraystretch}{1.25}
        \begin{tabular}{l c c c c c}
            \toprule
            \textbf{Escenario}           & \textbf{Biomasa (t/año)} & \textbf{Bioetanol (t/año)} & \textbf{Lignina Pura (t/año)} & \textbf{Biogás (\unit{\cubic\meter}/año)} & \textbf{CO$_2$ (t/año)} \\
            \midrule
            B-1 (Solo Hueso)             & \num{15000}              & \num{2028}                 & \num{2093}                    & \num{3912361}                             & \num{5904}              \\
            B-2 (Hoja de Olivo)          & \num{40000}              & \num{4166}                 & --                            & --                                        & --                      \\
            B-3 (Orujo Desalperujado)    & \num{40000}              & \num{5696}                 & --                            & --                                        & --                      \\
            B-4 (Madera de Poda)         & \num{65000}              & \num{9478}                 & --                            & --                                        & --                      \\
            B-5 (Mezcla Integral Olivar) & \num{160000}             & \num{23444}                & \num{14197}                   & \num{22460957}                            & \num{66134}             \\
            \bottomrule
        \end{tabular}%
    }
    \label{tab:escenarios_referencia}
\end{table}

\section{Criterios de Escala: Del Hueso a la Mezcla Integral de Residuos}
El diseño de la capacidad de la planta debe conciliar dos factores contrapuestos: la curva de economía de escala tecnoeconómica y el radio logístico óptimo de transporte de biomasa \autocite{Jafari-Gohargan2026}:
\begin{enumerate}
    \item \textbf{Capacidad Mínima Viable ($C_{\text{min}}$):} Según \textcite{Jafari-Gohargan2026}, el escenario B-1 (solo hueso de aceituna a \SI{15000}{\tonne/a\text{\~{n}}o}) resulta deficitario con un Índice de Rentabilidad ($PI$) de \num{0.80} y una tasa de retorno neta ($NRR$) de \SI{-19.98}{\percent}, situándose su capacidad crítica de viabilidad ($PI \ge 1.0$) en \SI{30000}{\tonne/a\text{\~{n}}o}.
    \item \textbf{Alimentación Multibiomasa (Hueso, Hoja y Madera de Poda):} Al ampliar el parque de materia prima hacia los subproductos agrícolas de la poda bianual del olivar (hoja y ramas/madera de poda), se incrementa drásticamente la resiliencia productiva y el aprovechamiento estacional:
          \begin{itemize}
              \item \textbf{Campaña de Molturación (Noviembre -- Febrero):} Recepción intensiva de hueso fresco desde almazaras (humedad \SI{8}{\percent} -- \SI{12}{\percent}).
              \item \textbf{Campaña de Poda (Enero -- Abril):} Trituración en campo de madera y follaje de poda, estabilizando el caudal de alimentación a la biorrefinería a lo largo de las \SI{8000}{\hour/a\text{\~{n}}o} de diseño.
          \end{itemize}
    \item \textbf{Escala Nominal Seleccionada:} Se establece un régimen nominal de \SI{45000}{\tonne/a\text{\~{n}}o} en base hueso de aceituna (equivalente a una planta modular ampliable hasta \SI{100000}{\tonne/a\text{\~{n}}o} en régimen multibiomasa B-5 con hoja y poda).
\end{enumerate}

\section{Balance Productivo e Ingresos Brutos Estimados}
Ejecutando el modelo implementado en \texttt{estudio\_mercado\_localizacion.py}, se obtienen los consumos de materias primas auxiliares, los flujos de producción anual y los ingresos brutos facturables para la escala base de \SI{45000}{\tonne/a\text{\~{n}}o} (Tabla~\ref{tab:balance_45k}).

\begin{table}[H]
    \centering
    \caption{Balance de producción y ventas anuales para una planta de \SI{45000}{\tonne/a\text{\~{n}}o}.}
    \adjustbox{max width=0.95\textwidth}{%
        \renewcommand{\arraystretch}{1.3}
        \begin{tabular}{l c c c}
            \toprule
            \textbf{Concepto / Corriente}                                & \textbf{Flujo Anual}                       & \textbf{Precio Ref. (\$/unidad)} & \textbf{Ingresos Anuales (\$/año)} \\
            \midrule
            \textbf{Bioetanol 2G (Corriente 40)}                         & \SI{6084.0}{\tonne/a\text{\~{n}}o}         & \SI{1200}{\$/\tonne}             & \$\num{7300800.00}                 \\
            \textbf{Lignina Pura Técnica (Corriente 19)}                 & \SI{6279.0}{\tonne/a\text{\~{n}}o}         & \SI{630}{\$/\tonne}              & \$\num{3955770.00}                 \\
            \textbf{Fracción Rica en Lignina (coproducto)}               & \SI{28509.0}{\tonne/a\text{\~{n}}o}        & \SI{42}{\$/\tonne}               & \$\num{1197378.00}                 \\
            \textbf{Dióxido de Carbono Biogénico (Corriente 30)}         & \SI{17712.0}{\tonne/a\text{\~{n}}o}        & \SI{38}{\$/\tonne}               & \$\num{673056.00}                  \\
            \textbf{Biofertilizante Deshidratado (Corriente 26)}         & \SI{1950.0}{\tonne/a\text{\~{n}}o}         & \SI{4.5}{\$/\tonne}              & \$\num{8775.00}                    \\
            \midrule
            \textbf{Biogás Generado (Corriente 27)}                      & \SI{11737083}{\cubic\meter/a\text{\~{n}}o} & Autoconsumo                      & Caldera Cogeneración               \\
            \midrule
            \multicolumn{3}{l}{\textbf{TOTAL INGRESOS BRUTOS ESTIMADOS}} & \textbf{\$\num{13135779.00}/año}                                                                                   \\
            \bottomrule
        \end{tabular}%
    }
    \label{tab:balance_45k}
\end{table}

\begin{tcolorbox}[colback=blue!5!white,colframe=blue!75!black,title=\textbf{Consumos de Reactivos Clave Calculados}]
    Para procesar las \SI{45000}{\tonne/a\text{\~{n}}o} de hueso de aceituna, el proceso requiere anualmente:
    \begin{itemize}
        \item Complejo Enzimático Cellic CTec2: \SI{210.0}{\kilo\gram/a\text{\~{n}}o}.
        \item Ácido Sulfúrico ($\text{H}_2\text{SO}_4$ \SI{98}{\percent}): \SI{411.0}{\tonne/a\text{\~{n}}o}.
        \item Amoniaco anhidro ($\text{NH}_3$): \SI{144.0}{\tonne/a\text{\~{n}}o}.
        \item Materia prima hueso de aceituna (\SI{17}{\$/\tonne}): \$\num{765000.00}/año.
    \end{itemize}
\end{tcolorbox}

%=============================================================================
\chapter{Estudio de Localización de la Planta Industrial}

\section{Justificación Cuantitativa de Localización con SABI y AICA}
Para fundamentar objetivamente la macrolocalización, se ha realizado una prospección comarcal en la base de datos empresarial \textbf{SABI (Informa)} cruzada con los aforos de \textbf{AICA} y el censo industrial de \textbf{AVEBIOM}:
\begin{enumerate}
    \item Se han extraído las empresas fabricantes de aceite de oliva bajo el código \textbf{CNAE 1043} en las provincias de Granada, Córdoba y Jaén.
    \item Se ha agregado la cifra de ingresos de explotación comarcal ($\sum \text{Ingresos CNAE 1043}$ en millones de euros), como indicador directo de la masa crítica de aceituna molturada y hueso producido.
    \item Se ha contabilizado la densidad de plantas de biomasa, orujeras extractoras y secaderos industriales que operan como competidores directos en la compra de hueso para quema térmica.
    \item Se calcula el ratio de oportunidad:
          \begin{equation}
              \text{Ratio SABI} = \frac{\sum \text{Ingresos de Explotación CNAE 1043 (M€)}}{\text{Número de Competidores de Biomasa y Orujeras}}
          \end{equation}
\end{enumerate}
La comarca con el mayor ratio representa la ubicación óptima: máxima abundancia de materia prima con menor canibalización de precios (Tabla~\ref{tab:sabi_ratio}).

\begin{table}[H]
    \centering
    \caption{Análisis del ratio de oportunidad en SABI por comarcas oleícolas.}
    \adjustbox{max width=0.98\textwidth}{%
        \renewcommand{\arraystretch}{1.3}
        \begin{tabular}{l l c c c c}
            \toprule
            \textbf{Comarca Oleícola}                             & \textbf{Provincia} & \textbf{Ingresos CNAE 1043} & \textbf{Almazaras} & \textbf{Competidores Biomasa} & \textbf{Ratio SABI (M€/comp.)} \\
            \midrule
            \textbf{Poniente Granadino (Loja / Montefrío)}        & \textbf{Granada}   & \textbf{385.4 M€}           & \num{42}           & \num{3}                       & \textbf{128.47 (Ganadora)}     \\
            \textbf{Montes de Granada (Íznalloz)}                 & Granada            & 215.2 M€                    & \num{28}           & \num{2}                       & 107.60                         \\
            \textbf{Subbética Cordobesa (Lucena / Priego)}        & Córdoba            & 590.8 M€                    & \num{68}           & \num{7}                       & 84.40                          \\
            \textbf{Campiña Sur Cordobesa (Puente Genil)}         & Córdoba            & 410.5 M€                    & \num{46}           & \num{5}                       & 82.10                          \\
            \textbf{Sierra Sur de Jaén (Alcalá la Real / Martos)} & Jaén               & 695.0 M€                    & \num{74}           & \num{11}                      & 63.18                          \\
            \textbf{Campiña de Jaén (Mengíbar / Andújar)}         & Jaén               & 780.2 M€                    & \num{85}           & \num{14}                      & 55.73                          \\
            \bottomrule
        \end{tabular}%
    }
    \label{tab:sabi_ratio}
\end{table}

\section{Macrolocalización: Método de Factores Ponderados}
A partir del análisis comarcal anterior, se formaliza la evaluación provincial multicriterio. Los factores se ponderan y puntúan según las fuentes oficiales citadas:
\begin{itemize}
    \item \textbf{Disponibilidad de materia prima ($w_1 = \num{0.25}$):} Jaén (\num{10.0}), Córdoba (\num{9.0}) y Granada (\num{8.0}), de acuerdo con los aforos del MAPA/AICA.
    \item \textbf{Baja competencia térmica / Ratio SABI ($w_2 = \num{0.20}$):} Granada (\num{9.5}) lidera de forma concluyente gracias a su ratio de \num{128.47} M€/competidor frente a Jaén (\num{4.0}), donde las orujeras saturan la compra del hueso a precios elevados \autocite{AVEBIOM2026}.
    \item \textbf{Infraestructura energética ($w_3 = \num{0.20}$):} Capacidad de evacuación eléctrica en AT y proximidad al gasoducto de transporte primario (Red Eléctrica / Enagás).
    \item \textbf{Cercanía a clientes / Off-takers ($w_4 = \num{0.15}$):} Proximidad a centros logísticos de hidrocarburos y refinerías (Moeve San Roque, Huelva y CLH Cartagena) \autocite{CORES2026}.
    \item \textbf{Comunicaciones viarias ($w_5 = \num{0.10}$):} Enlace directo con los corredores A-92 (Transversal de Andalucía) y A-44 (Autovía de Sierra Nevada).
    \item \textbf{Recursos hídricos y vertido ($w_6 = \num{0.10}$):} Disponibilidad de dotación industrial en cuenca regulada (Confederación Hidrográfica del Guadalquivir / Cuenca Mediterránea).
\end{itemize}

\begin{table}[H]
    \centering
    \caption{Matriz de Macrolocalización por Factores Ponderados.}
    \adjustbox{max width=0.98\textwidth}{%
        \renewcommand{\arraystretch}{1.3}
        \begin{tabular}{l c c c c}
            \toprule
            \textbf{Criterio / Factor de Decisión}                  & \textbf{Peso ($w_i$)} & \textbf{Jaén} & \textbf{Córdoba} & \textbf{Granada} \\
            \midrule
            Disponibilidad materia prima (Hueso + Poda - MAPA/AICA) & \num{0.25}            & \num{10.0}    & \num{9.0}        & \num{8.0}        \\
            Baja competencia térmica / Ratio SABI (AVEBIOM/SABI)    & \num{0.20}            & \num{4.0}     & \num{7.0}        & \num{9.5}        \\
            Infraestructura energética (Red AT y Gasoducto)         & \num{0.20}            & \num{8.0}     & \num{8.0}        & \num{9.0}        \\
            Cercanía a clientes (Off-takers etanol CORES/Moeve)     & \num{0.15}            & \num{7.0}     & \num{9.0}        & \num{8.0}        \\
            Comunicaciones y logística viaria (A-92/A-44/A-45)      & \num{0.10}            & \num{8.0}     & \num{9.0}        & \num{9.0}        \\
            Disponibilidad de agua y gestión de efluentes (CHG)     & \num{0.10}            & \num{7.0}     & \num{7.0}        & \num{8.0}        \\
            \midrule
            \textbf{PUNTUACIÓN TOTAL PONDERADA}                     & \textbf{1.00}         & \textbf{7.45} & \textbf{8.20}    & \textbf{8.60}    \\
            \bottomrule
        \end{tabular}%
    }
    \label{tab:macrolocalizacion}
\end{table}

\begin{figure}[H]
    \centering
    \begin{subfigure}[b]{0.48\textwidth}
        \centering
        \adjustbox{max width=\linewidth}{%
            \begin{tikzpicture}[scale=0.72]
                % Malla poligonal de fondo regular (escala 0 a 10)
                \foreach \r in {2, 4, 6, 8, 10} {
                        \draw[gray!35, thin] (90:0.35*\r) -- (30:0.35*\r) -- (330:0.35*\r) -- (270:0.35*\r) -- (210:0.35*\r) -- (150:0.35*\r) -- cycle;
                        \node[gray!65, font=\tiny, anchor=south west] at (90:0.35*\r) {\r};
                    }

                % Ejes radiales
                \foreach \ang in {90, 30, 330, 270, 210, 150} {
                        \draw[gray!50, dashed] (0,0) -- (\ang:3.8);
                    }

                % Etiquetas de criterios
                \node[above, font=\footnotesize\bfseries, align=center] at (90:3.85) {Disponibilidad MP\\(Peso 25\,\%)};
                \node[above right, font=\footnotesize\bfseries, align=left] at (30:3.85) {Baja Comp. Térmica\\(Ratio SABI, 20\,\%)};
                \node[below right, font=\footnotesize\bfseries, align=left] at (330:3.85) {Infraestructura\\Energética (20\,\%)};
                \node[below, font=\footnotesize\bfseries, align=center] at (270:3.85) {Cercanía Clientes\\(Off-takers, 15\,\%)};
                \node[below left, font=\footnotesize\bfseries, align=right] at (210:3.85) {Comunicaciones\\Viarias (10\,\%)};
                \node[above left, font=\footnotesize\bfseries, align=right] at (150:3.85) {Recursos\\Hídricos (10\,\%)};

                % Polígonos de puntuación directa (radio = 0.35 * nota)
                % Granada: [8, 9.5, 9, 8, 9, 8]
                \fill[green!70!black, fill opacity=0.20]
                (90:2.80) -- (30:3.325) -- (330:3.15) -- (270:2.80) -- (210:3.15) -- (150:2.80) -- cycle;
                \draw[green!60!black, line width=1.4pt, mark=triangle*, mark size=2pt, mark options={fill=green!70!black}]
                (90:2.80) -- (30:3.325) -- (330:3.15) -- (270:2.80) -- (210:3.15) -- (150:2.80) -- cycle;

                % Córdoba: [9, 7, 8, 9, 9, 7]
                \fill[orange!85!red, fill opacity=0.15]
                (90:3.15) -- (30:2.45) -- (330:2.80) -- (270:3.15) -- (210:3.15) -- (150:2.45) -- cycle;
                \draw[orange!90!black, very thick, dashed, mark=square*, mark size=1.8pt, mark options={fill=orange!80}]
                (90:3.15) -- (30:2.45) -- (330:2.80) -- (270:3.15) -- (210:3.15) -- (150:2.45) -- cycle;

                % Jaén: [10, 4, 8, 7, 8, 7]
                \fill[blue!70!cyan, fill opacity=0.15]
                (90:3.50) -- (30:1.40) -- (330:2.80) -- (270:2.45) -- (210:2.80) -- (150:2.45) -- cycle;
                \draw[blue!85!black, very thick, mark=*, mark size=1.8pt, mark options={fill=blue!80}]
                (90:3.50) -- (30:1.40) -- (330:2.80) -- (270:2.45) -- (210:2.80) -- (150:2.45) -- cycle;
            \end{tikzpicture}%
        }
        \caption{Puntuación directa intrínseca (escala 0--10).}
        \label{fig:radar_directo}
    \end{subfigure}\hfill
    \begin{subfigure}[b]{0.48\textwidth}
        \centering
        \adjustbox{max width=\linewidth}{%
            \begin{tikzpicture}[scale=0.72]
                % Malla proporcional a los pesos máximos (w_i * 10)
                % Max radios: 90: 3.5cm (2.5), 30: 2.8cm (2.0), 330: 2.8cm (2.0), 270: 2.1cm (1.5), 210: 1.4cm (1.0), 150: 1.4cm (1.0)
                % Factor de escala: 1.4 cm por unidad ponderada (2.5 * 1.4 = 3.5cm)
                \foreach \frac in {0.25, 0.50, 0.75, 1.0} {
                        \draw[gray!35, thin]
                        (90:3.5*\frac) -- (30:2.8*\frac) -- (330:2.8*\frac) -- (270:2.1*\frac) -- (210:1.4*\frac) -- (150:1.4*\frac) -- cycle;
                    }
                \node[gray!65, font=\tiny, anchor=south west] at (90:3.5) {2.50};
                \node[gray!65, font=\tiny, anchor=south west] at (30:2.8) {2.00};
                \node[gray!65, font=\tiny, anchor=north west] at (270:2.1) {1.50};
                \node[gray!65, font=\tiny, anchor=north east] at (210:1.4) {1.00};

                % Ejes radiales con longitud igual a su contribución máxima
                \draw[gray!50, dashed] (0,0) -- (90:3.8);
                \draw[gray!50, dashed] (0,0) -- (30:3.1);
                \draw[gray!50, dashed] (0,0) -- (330:3.1);
                \draw[gray!50, dashed] (0,0) -- (270:2.4);
                \draw[gray!50, dashed] (0,0) -- (210:1.7);
                \draw[gray!50, dashed] (0,0) -- (150:1.7);

                % Etiquetas de criterios con peso máximo
                \node[above, font=\footnotesize\bfseries, align=center] at (90:3.85) {Disponibilidad MP\\(Máx. 2.50 pts)};
                \node[above right, font=\footnotesize\bfseries, align=left] at (30:3.15) {Baja Comp. Térmica\\(Máx. 2.00 pts)};
                \node[below right, font=\footnotesize\bfseries, align=left] at (330:3.15) {Infraestructura\\(Máx. 2.00 pts)};
                \node[below, font=\footnotesize\bfseries, align=center] at (270:2.45) {Clientes Off-takers\\(Máx. 1.50 pts)};
                \node[below left, font=\footnotesize\bfseries, align=right] at (210:1.75) {Comunicaciones\\(Máx. 1.00 pts)};
                \node[above left, font=\footnotesize\bfseries, align=right] at (150:1.75) {Agua y Vertido\\(Máx. 1.00 pts)};

                % Polígonos de contribución ponderada real: radio = 1.4 * (w_i * P_ij)
                % Granada: [2.00, 1.90, 1.80, 1.20, 0.90, 0.80] -> radio = 1.4 * contrib
                % R_G = [2.80, 2.66, 2.52, 1.68, 1.26, 1.12]
                \fill[green!70!black, fill opacity=0.20]
                (90:2.80) -- (30:2.66) -- (330:2.52) -- (270:1.68) -- (210:1.26) -- (150:1.12) -- cycle;
                \draw[green!60!black, line width=1.4pt, mark=triangle*, mark size=2pt, mark options={fill=green!70!black}]
                (90:2.80) -- (30:2.66) -- (330:2.52) -- (270:1.68) -- (210:1.26) -- (150:1.12) -- cycle;

                % Córdoba: [2.25, 1.40, 1.60, 1.35, 0.90, 0.70] -> radio = 1.4 * contrib
                % R_C = [3.15, 1.96, 2.24, 1.89, 1.26, 0.98]
                \fill[orange!85!red, fill opacity=0.15]
                (90:3.15) -- (30:1.96) -- (330:2.24) -- (270:1.89) -- (210:1.26) -- (150:0.98) -- cycle;
                \draw[orange!90!black, very thick, dashed, mark=square*, mark size=1.8pt, mark options={fill=orange!80}]
                (90:3.15) -- (30:1.96) -- (330:2.24) -- (270:1.89) -- (210:1.26) -- (150:0.98) -- cycle;

                % Jaén: [2.50, 0.80, 1.60, 1.05, 0.80, 0.70] -> radio = 1.4 * contrib
                % R_J = [3.50, 1.12, 2.24, 1.47, 1.12, 0.98]
                \fill[blue!70!cyan, fill opacity=0.15]
                (90:3.50) -- (30:1.12) -- (330:2.24) -- (270:1.47) -- (210:1.12) -- (150:0.98) -- cycle;
                \draw[blue!85!black, very thick, mark=*, mark size=1.8pt, mark options={fill=blue!80}]
                (90:3.50) -- (30:1.12) -- (330:2.24) -- (270:1.47) -- (210:1.12) -- (150:0.98) -- cycle;
            \end{tikzpicture}%
        }
        \caption{Contribución ponderada real a la nota final.}
        \label{fig:radar_ponderado}
    \end{subfigure}

    \vspace{0.8em}
    % Leyenda unificada inferior
    \begin{tikzpicture}
        \draw[gray!35, fill=white, rounded corners=3pt] (-0.2,-0.35) rectangle (10.2,0.45);
        \draw[green!60!black, line width=1.5pt, mark=triangle*, mark size=2.5pt, mark options={fill=green!70!black}] (0.2,0.05) -- (0.8,0.05);
        \node[font=\footnotesize\bfseries, anchor=west, text=green!50!black] at (0.9, 0.05) {Granada (\num{8.60}/10)};

        \draw[orange!90!black, very thick, dashed, mark=square*, mark size=2pt, mark options={fill=orange!80}] (3.7,0.05) -- (4.3,0.05);
        \node[font=\footnotesize, anchor=west] at (4.4, 0.05) {Córdoba (\num{8.20}/10)};

        \draw[blue!85!black, very thick, mark=*, mark size=2pt, mark options={fill=blue!80}] (7.2,0.05) -- (7.8,0.05);
        \node[font=\footnotesize, anchor=west] at (7.9, 0.05) {Jaén (\num{7.45}/10)};
    \end{tikzpicture}
    \caption{Comparativa multicriterio de macrolocalización: (a) Desempeño técnico intrínseco y (b) Impacto ponderado en la función objetivo.}
    \label{fig:radar_macrolocalizacion}
\end{figure}

La provincia de **Granada obtiene la mayor calificación (8.60/10)**, impulsada por el Poniente Granadino como la zona con mayor ratio de excedente de biomasa respecto a competidores (Figura~\ref{fig:radar_macrolocalizacion}).

\section{Microlocalización: Método de Brown y Gibson}
Para la selección del emplazamiento exacto se aplica el algoritmo cuantitativo de **Brown y Gibson**, combinando factores críticos ($FC_i \in \{0, 1\}$), factores objetivos basados en costes anuales ($FO_i$) y factores subjetivos ponderados ($FS_i$):
\begin{equation}
    MPL_i = FC_i \cdot \left[ \alpha \cdot FO_i + (1 - \alpha) \cdot FS_i \right]
\end{equation}
donde se adopta un peso objetivo $\alpha = \num{0.70}$ y una ponderación subjetiva $1 - \alpha = \num{0.30}$.
Los factores objetivos computan la suma anualizada de:
\begin{itemize}
    \item Coste de adquisición y urbanización de suelo industrial.
    \item Coste logístico de transporte de materia prima y distribución de productos.
    \item Coste de acometidas eléctricas, gas natural y canon de vertidos.
\end{itemize}

\begin{table}[H]
    \centering
    \caption{Resultados del Método de Brown y Gibson para microlocalización.}
    \adjustbox{max width=0.98\textwidth}{%
        \renewcommand{\arraystretch}{1.35}
        \begin{tabular}{l c c c c c}
            \toprule
            \textbf{Polígono Industrial Candidato}                   & \textbf{$FC_i$} & \textbf{Coste Anual (€)} & \textbf{$FO_i$ (Objetivo)} & \textbf{$FS_i$ (Subjetivo)} & \textbf{$MPL_i$ (Índice Final)} \\
            \midrule
            \textbf{Polígono Manzanil II / Loja (Granada)}           & \num{1.0}       & \num{1450000}            & \num{0.3503}               & \num{0.3800}                & \textbf{0.3592 (1º)}            \\
            \textbf{Polígono Los Santos / Lucena (Córdoba)}          & \num{1.0}       & \num{1520000}            & \num{0.3342}               & \num{0.3500}                & \textbf{0.3389 (2º)}            \\
            \textbf{Polígono Llano Mazuelos / Alcalá la Real (Jaén)} & \num{1.0}       & \num{1610000}            & \num{0.3155}               & \num{0.2700}                & \textbf{0.3019 (3º)}            \\
            \bottomrule
        \end{tabular}%
    }
    \label{tab:brown_gibson}
\end{table}

\begin{tcolorbox}[colback=green!5!white,colframe=green!60!black,title=\textbf{Emplazamiento Industrial Seleccionado}]
    El emplazamiento ganador indiscutible es el **Polígono Industrial Manzanil II en Loja (Granada)**, con un índice $MPL = \mathbf{0.3592}$. Destaca por su menor coste operativo global (\num{1450000} €/año), conexión directa con la autovía A-92 y disponibilidad inmediata de parcelas industriales aptas para instalaciones químicas catalogadas.
\end{tcolorbox}

%=============================================================================
\chapter{Esquema General de Procesos e Integración}

\section{Diagrama de Bloques y Flujo del Proceso}
La planta integra secuencialmente las operaciones de acondicionamiento físico, doble hidrólisis, fermentaciones diferenciadas y valorización integral de efluentes (Figura~\ref{fig:diagrama_bloques_esquema}).

\begin{figure}[H]
    \centering
    \adjustbox{max width=0.98\textwidth}{%
        \begin{tikzpicture}[>=Latex, node distance=1.6cm and 1.8cm, auto,
                block/.style={rectangle, draw=blue!80!black, fill=blue!10, rounded corners=4pt, thick, text width=3.3cm, align=center, minimum height=1.35cm, font=\small\bfseries},
                block_sep/.style={rectangle, draw=orange!80!black, fill=orange!15, rounded corners=4pt, thick, text width=3.3cm, align=center, minimum height=1.35cm, font=\small\bfseries},
                block_rec/.style={rectangle, draw=green!70!black, fill=green!15, rounded corners=4pt, thick, text width=3.3cm, align=center, minimum height=1.35cm, font=\small\bfseries},
                line/.style={draw=black!85, line width=1.2pt, ->},
                line_ret/.style={draw=green!60!black, line width=1.2pt, dashed, ->}]

            % Fila 1: Acondicionamiento y Fraccionamiento Secuencial
            \node [block] (prep) {Acondicionamiento\\(Secado, Triturado y Tamizado)};
            \node [block, right=1.8cm of prep] (hidro1) {Etapa 1: Pretratamiento\\Ácido Diluido ($\text{H}_2\text{SO}_4$)};
            \node [block_sep, right=1.8cm of hidro1] (organo) {Etapa 2: Reactor\\Organosolv (Etanol/Agua)};

            % Fila 2: Fermentación Pentosas y Hexosas
            \node [block, below=1.6cm of hidro1] (ferm_pent) {Fermentación Pentosas C5\\(Cepa \textit{E. coli})};
            \node [block, below=1.6cm of organo] (sacar) {Presacarificación y SSF C6\\(Cepa \textit{S. cerevisiae})};
            \node [block_rec, right=1.8cm of organo] (lignina) {Precipitación y Secado\\de Lignina Pura};

            % Fila 3: Separación/Destilación y Digestión Anaerobia
            \node [block_sep, below=1.6cm of prep] (destil) {Separación y Destilación\\Etanol 2G / Furfural};
            \node [block_rec, below=1.6cm of ferm_pent] (anaerobio) {Digestión Anaerobia\\y Cogeneración Biogás};

            % Flechas principales
            \path [line] (prep) -- node[above, font=\footnotesize] {Hueso seco} (hidro1);
            \path [line] (hidro1) -- node[above, font=\footnotesize] {Fracción sólida} (organo);
            \path [line] (organo) -- node[above, font=\footnotesize] {Licor negro} (lignina);

            \path [line] (hidro1) -- node[right, font=\footnotesize] {Licor C5} (ferm_pent);
            \path [line] (organo) -- node[right, font=\footnotesize] {Fracción C6} (sacar);

            % Corrientes de caldo fermentado hacia Destilación
            \path [line] (ferm_pent) -- node[above, font=\footnotesize] {Mosto C5} (destil);
            \path [line] (sacar.south) -- ++(0,-0.45) -| node[pos=0.25, below, font=\footnotesize] {Mosto fermentado C6} (destil.south);

            % Vinazas residuales de fondo de destilación a Digestor Anaerobio
            \path [line] (destil.east) -- node[above, font=\footnotesize] {Vinazas residuales} (anaerobio.west);

            % Retorno de Biogás perimetral limpio hasta Acondicionamiento (Secado)
            \path [line_ret] (anaerobio.south) -- ++(0,-0.7) -| ([xshift=-0.9cm]prep.west) -- node[pos=0.8, left, font=\footnotesize, text=green!50!black] {\shortstack{Biometano\\(Secador rotatorio)}} (prep.west);

        \end{tikzpicture}%
    }
    \caption{Diagrama de bloques general e integración energética de la biorrefinería.}
    \label{fig:diagrama_bloques_esquema}
\end{figure}

\section{Síntesis de Equipos e Identificación en Diagramas}
Siguiendo los planos de ingeniería generados en el proyecto (\texttt{D-bloques.pdf} y \texttt{D-flujo.pdf}):
\begin{itemize}
    \item **Línea de Acondicionamiento:** Cinta transportadora (\texttt{CT-01}), Secador rotatorio (\texttt{S-01}), Cinta (\texttt{CT-02}), Molino triturador de martillos (\texttt{T-01}) y Tamiz clasificador rotativo (\texttt{TR-01}) con reciclo de gruesos (corriente 6).
    \item **Línea de Pentosas:** Reactor tubular continuo de hidrólisis ácida (\texttt{R-01}), Filtro prensa (\texttt{F-01}), Columna de adsorción con carbón activo (\texttt{A-01}) y Biorreactor continuo de pentosas (\texttt{BF-01}).
    \item **Línea de Hexosas y Lignina:** Reactor presurizado Organosolv (\texttt{R-02}), Cámara Flash de despresurización (\texttt{V-101}), Filtro de banda (\texttt{F-02}), Reactor de presacarificación enzimática (\texttt{R-03}) y Biorreactores SSF (\texttt{BF-02}).
    \item **Tren de Destilación y Cogeneración:** Columna fraccionadora principal (\texttt{Z-01}), Divisor de reflujo (\texttt{D-01}), Mezclador de reciclo hidroalcohólico (\texttt{M-01}), Digestor anaerobio de lodos (\texttt{DA-01}) y Caldera de vapor de cogeneración (\texttt{Q-01}).
\end{itemize}

%=============================================================================
\printbibliography[heading=bibintoc, title={Referencias Bibliográficas}]

\end{document}
